\clearpage\selectlanguage{german}\fancyhead[L]{\small ARITHMETISCH GEOMETRISCHES MITTEL.}
\pdfbookmark[0]{FORTSETZUNG DER UNTERSUCHUNGEN ÜBER DAS ARITHMETISCH GEOMETRISCHE MITTEL}{agm-fortsetzung}
\SourceLeaf{384}
\begin{center}\bfseries
FORTSETZUNG DER UNTERSUCHUNGEN

ÜBER DAS ARITHMETISCH GEOMETRISCHE MITTEL.

\rule{0.08\linewidth}{0.3pt}
\end{center}
[Die weiteren Untersuchungen von \textsc{Gauss} über das Arithmetisch geometrische Mittel sind so unvollständig, nur in den Endformeln, und mit so wechselnder Bezeichnung niedergeschrieben, dass fast alle Übersicht fehlen würde, wenn der Abdruck nur jene Formeln ohne nebenhergehende Erläuterung wiedergäbe. In den folgenden Artikeln habe ich die im Nachlasse gefundenen Stellen hervorgehoben und durch solchen Gedankengang verbunden, wie er vielleicht die Veranlassung gewesen ist, dass \textsc{Gauss} die betreffenden Resultate unter einem erkennbaren gemeinsamen Gesichtspunkte aufgestellt hat.]

\AGMArticle{12}
[An mehren Stellen bedient sich \textsc{Gauss} neben den in den vorhergehenden Artikeln betrachteten beiden ursprünglichen Reihen von Grössen $(a,b)$, welche sich auf dasselbe arithmetisch geometrische Mittel beziehen, auch noch einer dritten Reihe von Grössen $(c)$, die mit jenen durch die Gleichung $aa=bb+cc$, welche für jeden gemeinsamen Index der drei Glieder $(a,b,c)$ gilt, zusammenhängt. Nach den Festsetzungen in Art. 2 können die Gleichungen zur Berechnung der rück- und vorwärts folgenden aus $a,b,c$ abgeleiteten Grössen in die Form gebracht werden:
\SourceLeaf{385}
\[
\begin{aligned}&.\ .\quad{}^{\prime\prime}a={}^\prime a+{}^\prime c,\quad{}^\prime a=a+c,\quad a,\\&\qquad2a'=a+b,\quad2a''=a'+b',\quad.\ .\\[9pt]&.\ .\quad{}^{\prime\prime}b={}^\prime a-{}^\prime c,\quad{}^\prime b=a-c,\quad b,\\&\qquad b'b'=ab; \quad b''b''=a'b',\quad.\ .\\[9pt]&.\ .\quad{}^{\prime\prime}c{}^{\prime\prime}c=4{}^\prime a{}^\prime c,\quad{}^\prime c{}^\prime c=4ac,\quad c,\\&\qquad2c'=a-b,\quad2c''=a'-b',\quad.\ .\end{aligned}
\]
Hieraus ergibt sich unmittelbar dass, während
\[
\begin{aligned}&\mathrm{M}(a,b)=\mathrm{M}(a^{\mathrm{n}},b^{\mathrm{n}})=\mathrm{M}({}^{\mathrm{n}}a,{}^{\mathrm{n}}b)\\&=\lim a^{\mathrm{m}}=\lim b^{\mathrm{m}}\quad\hbox{für}\quad m=\infty\end{aligned}
\]
ist,
\[
\begin{aligned}&\mathrm{M}(a,c)=2^n\mathrm{M}(a^{\mathrm{n}},c^{\mathrm{n}})=2^{-n}\mathrm{M}({}^{\mathrm{n}}a,{}^{\mathrm{n}}c)\\&=\lim\frac{{}^{\mathrm{m}}a}{2^m}=\lim\frac{{}^{\mathrm{m}}c}{2^m}\quad\hbox{für}\quad m=\infty\end{aligned}
\]
wird und sich damit die dem ersten Beispiel in Art. 3 angefügte Bemerkung, dass die ${}^{\mathrm{n}}a,{}^{\mathrm{n}+1}a,{}^{\mathrm{n}+2}a$ etwa wie die Glieder einer geometrischen Reihe mit dem Quotienten $2$ zunehmen als allgemein gültig erweist, ferner dass für den im Beispiel 4 Art. 3 berechneten Fall $a=b\surd2=c\surd2$ bei jedem $n$: ${}^{\mathrm{n}}a=2^n.a^{\mathrm{n}}$, ${}^{\mathrm{n}}b=2^n.c^{\mathrm{n}}$, ${}^{\mathrm{n}}c=2^n.b^{\mathrm{n}}$ ist.

\par Die Art der Annäherung jener Grössen an ihre Grenzwerthe ergibt sich aus den folgenden Reihenentwicklungen
\[
\begin{aligned}\mathrm{M}(a,b)&=a^{\mathrm{n}}-c^{\mathrm{n}+1}-c^{\mathrm{n}+2}\\&\quad-.\;.\;.-c^{\mathrm{n}+\mathrm{m}}-.\;.\;.\end{aligned}
\]
\[
\begin{aligned}\mathrm{M}(a,b)&=b^{\mathrm{n}}+c^{\mathrm{n}+1}-c^{\mathrm{n}+2}\\&\quad-.\;.\;.-c^{\mathrm{n}+\mathrm{m}}-.\;.\;.\end{aligned}
\]
\[
\begin{aligned}\mathrm{M}(a,c)&=2^{-n}.{}^{\mathrm{n}}a-2^{-n-1}.{}^{\mathrm{n}+1}b-2^{-n-2}.{}^{\mathrm{n}+2}b\\&\quad-.\;.\;.-2^{-n-m}.{}^{\mathrm{n}+\mathrm{m}}b-.\;.\;.\end{aligned}
\]
\[
\begin{aligned}\mathrm{M}(a,c)&=2^{-n}.{}^{\mathrm{n}}c+2^{-n-1}.{}^{\mathrm{n}+1}b-2^{-n-2}.{}^{\mathrm{n}+2}b\\&\quad-.\;.\;.-2^{-n-m}.{}^{\mathrm{n}+\mathrm{m}}b-.\;.\;.\end{aligned}
\]
\[
\begin{aligned}\mathrm{M}(a,b)^2&=a^{\mathrm{n}}.a^{\mathrm{n}}-\frac12c^{\mathrm{n}}.c^{\mathrm{n}}-\frac32c^{\mathrm{n}+1}.c^{\mathrm{n}+1}\\&\quad-.\;.\;.-\frac32c^{\mathrm{n}+\mathrm{m}}.c^{\mathrm{n}+\mathrm{m}}-.\;.\;.\end{aligned}
\]
\[
\begin{aligned}\mathrm{M}(a,b)^2&=b^{\mathrm{n}}.b^{\mathrm{n}}+\frac12c^{\mathrm{n}}.c^{\mathrm{n}}-\frac32c^{\mathrm{n}+1}.c^{\mathrm{n}+1}\\&\quad-.\;.\;.-\frac32c^{\mathrm{n}+\mathrm{m}}.c^{\mathrm{n}+\mathrm{m}}-.\;.\;.\end{aligned}
\]
\[
\begin{aligned}\mathrm{M}(a,c)^2&=2^{-2n}.{}^{\mathrm{n}}c.{}^{\mathrm{n}}c+\frac12 2^{-2n}.{}^{\mathrm{n}}b.{}^{\mathrm{n}}b-\frac32 2^{-2n-2}.{}^{\mathrm{n}+1}b.{}^{\mathrm{n}+1}b\\&\quad-.\;.\;.\qquad-\frac32 2^{-2n-2m}.{}^{\mathrm{n}+\mathrm{m}}b.{}^{\mathrm{n}+\mathrm{m}}b-.\;.\;.\end{aligned}
\]
\[
\begin{aligned}\mathrm{M}(a,c)^2&=2^{-2n}.{}^{\mathrm{n}}a.{}^{\mathrm{n}}a-\frac12 2^{-2n}.{}^{\mathrm{n}}b.{}^{\mathrm{n}}b-\frac32 2^{-2n-2}.{}^{\mathrm{n}+1}b.{}^{\mathrm{n}+1}b\\&\quad-.\;.\;.\qquad-\frac32 2^{-2n-2m}.{}^{\mathrm{n}+\mathrm{m}}b.{}^{\mathrm{n}+\mathrm{m}}b-.\;.\;.\end{aligned}
\]
\[
\begin{aligned}\surd\mathrm{M}(a,b)&=\surd a^{\mathrm{n}}-\surd c^{\mathrm{n}+2}-\surd c^{\mathrm{n}+4}\\&\quad-.\;.\;.-\surd c^{\mathrm{n}+2\mathrm{m}}-.\;.\;.\end{aligned}
\]
\[
\begin{aligned}\surd\mathrm{M}(a,b)&=\surd b^{\mathrm{n}}+\surd c^{\mathrm{n}+2}-\surd c^{\mathrm{n}+4}\\&\quad-.\;.\;.-\surd c^{\mathrm{n}+2\mathrm{m}}-.\;.\;.\end{aligned}
\]
\SourceLeaf{386}
\[
\begin{aligned}\surd\mathrm{M}(a,c)&=\surd2^{-n}.{}^{\mathrm{n}}a-\surd2^{-n-2}.{}^{\mathrm{n}+2}b-\surd2^{-n-4}.{}^{\mathrm{n}+4}b\\&\quad-.\;.\;.-\surd2^{-n-2m}.{}^{\mathrm{n}+2\mathrm{m}}b-.\;.\;.\end{aligned}
\]
\[
\begin{aligned}\surd\mathrm{M}(a,c)&=\surd2^{-n}.{}^{\mathrm{n}}c+\surd2^{-n-2}.{}^{\mathrm{n}+2}b-\surd2^{-n-4}.{}^{\mathrm{n}+4}b\\&\quad-.\;.\;.-\surd2^{-n-2m}.{}^{\mathrm{n}+2\mathrm{m}}b-.\;.\;.\end{aligned}
\]
\[
\begin{aligned}\displaystyle\frac\pi2.\frac{\mathrm{M}(a,b)}{\mathrm{M}(a,c)}&\displaystyle=\frac1{2^n}\log4\frac{a^{\mathrm{n}}}{c^{\mathrm{n}}}-\frac1{2^n}\log\frac{a^{\mathrm{n}}}{a^{\mathrm{n}+1}}\\&\quad-.\;.\;.-\frac1{2^{n+m}}\log\frac{a^{\mathrm{n}+\mathrm{m}}}{a^{\mathrm{n}+\mathrm{m}+1}}-.\;.\end{aligned}
\]
\[
\begin{aligned}\displaystyle\frac\pi2.\frac{\mathrm{M}(a,b)}{\mathrm{M}(a,c)}&\displaystyle=\frac1{2^n}\log4\frac{b^{\mathrm{n}+1}}{c^{\mathrm{n}}}+\frac3{2^{n+1}}\log\frac{b^{\mathrm{n}+2}}{b^{\mathrm{n}+1}}\\&\quad+.\;.\;.+\frac3{2^{n+m}}\log\frac{b^{\mathrm{n}+\mathrm{m}+1}}{b^{\mathrm{n}+\mathrm{m}}}+.\;.\end{aligned}
\]
\[
\begin{aligned}\displaystyle\frac\pi2.\frac{\mathrm{M}(a,c)}{\mathrm{M}(a,b)}&\displaystyle=\frac1{2^n}\log2\frac{{}^{\mathrm{n}}a}{{}^{\mathrm{n}}b}-\frac1{2^n}\log2\frac{{}^{\mathrm{n}}a}{{}^{\mathrm{n}+1}a}\\&\quad-.\;.\;.-\frac1{2^{n+m}}\log2\frac{{}^{\mathrm{n}+\mathrm{m}}a}{{}^{\mathrm{n}+\mathrm{m}+1}a}-.\;.\end{aligned}
\]
\[
\begin{aligned}\displaystyle\frac\pi2.\frac{\mathrm{M}(a,c)}{\mathrm{M}(a,b)}&\displaystyle=\frac1{2^n}\log2\frac{{}^{\mathrm{n}+1}c}{{}^{\mathrm{n}}b}+\frac3{2^{n+1}}\log\frac12\frac{{}^{\mathrm{n}+2}c}{{}^{\mathrm{n}+1}c}\\&\quad+.\;.\;.+\frac3{2^{n+m}}\log\frac12\frac{{}^{\mathrm{n}+\mathrm{m}+1}c}{{}^{\mathrm{n}+\mathrm{m}}c}+.\;.\end{aligned}
\]
\par In den vier letzten Gleichungen ist $\frac\pi2$ statt der für beständig wachsendes $m$ geltenden Grenzwerthe beziehungsweise von
\[
\begin{gathered}\frac{\mathrm{M}(a^{\mathrm{m}},c^{\mathrm{m}})}{\mathrm{M}(a^{\mathrm{m}},b^{\mathrm{m}})}.\log4\frac{a^{\mathrm{m}}}{c^{\mathrm{m}}};\\[9pt]
\frac{\mathrm{M}(a^{\mathrm{m}},c^{\mathrm{m}})}{\mathrm{M}(a^{\mathrm{m}},b^{\mathrm{m}})}.\log4\frac{b^{\mathrm{m}+1}}{c^{\mathrm{m}}},\\[9pt]
\frac{\mathrm{M}({}^{\mathrm{m}}a,{}^{\mathrm{m}}b)}{\mathrm{M}({}^{\mathrm{m}}a,{}^{\mathrm{m}}c)}.\log4\frac{{}^{\mathrm{m}}a}{{}^{\mathrm{m}}b},\\[9pt]
\frac{\mathrm{M}({}^{\mathrm{m}}a,{}^{\mathrm{m}}b)}{\mathrm{M}({}^{\mathrm{m}}a,{}^{\mathrm{m}}c)}.\log2\frac{{}^{\mathrm{m}+1}c}{{}^{\mathrm{m}}b}\end{gathered}
\]
das ist statt des Grenzwerthes von $\mathrm{M}(1,\varepsilon)\log\frac4\varepsilon$ für bis zur Null abnehmende positive Werthe des $\varepsilon$ gesetzt. Es folgt nämlich zunächst aus jenen Gleichungen, in welchen, wie die Relationen
\[
\begin{gathered}\frac{a'}a=1-\frac{c'}a,\\[7pt]
\left(\frac{b'}b\right)^4=1+\frac{cc}{bb},\\[7pt]
\frac{{}^\prime a}a=2-\frac{{}^\prime b}a,\\[7pt]
\left(\frac12\frac{{}^\prime c}c\right)^4=1+\frac{bb}{cc}\end{gathered}
\]
leicht erkennen lassen, die zu logarithmirenden Werthe, so bald sie alle reell sind, vom zweiten Gliede der Reihe an beständig bis zur Einheit hin abnehmen, dass die gesuchte Grösse eine völlig bestimmte und z. B., wenn man $a=b\surd2=c\surd2$ setzt, eine zwischen den Grenzen $\frac94\log2$ und $\frac{10}4\log2$ eingeschlossene Grösse ist. Die obigen Gleichungen gelten auch für complexe Werthe der Veränderlichen, wenn man $n=0$ setzt und diejenigen Werthe der $\log{}^{\mathrm{m}}a$, $\log{}^{\mathrm{m}}c$, $\log a$, $\log b$, $\log c$, $\log a^{\mathrm{m}}$, $\log b^{\mathrm{m}}$ zu Grunde legt, welche durch stetige Änderung derselben aus den reellen Grössen folgen. Nimmt man nun als ein System $(a,b,c)$ das der Bedingung $a=b\surd2=c\surd2$ unterworfene und aus reellen Grössen bestehende, ferner als anderes System $(\alpha,\beta,\gamma)$ dasjenige, für welches $\alpha=c$, $\beta=b\surd{-1}$, $\gamma=a=b\surd2=c\surd2$ ist und die durch Wurzel-Ausziehung zu bestimmenden Werthe von $\beta^{\mathrm{m}}$ und ${}^{\mathrm{m}}\gamma$ positive reelle Theile erhalten, so kann man ein drittes veränderliches System $(A,B,C)$ aufstellen, welches von dem einen $(a,b,c)$ zu dem andern $(\alpha,\beta,\gamma)$ stetig übergeht, und zwar
\SourceLeaf{387}
so, dass bei diesem Übergang keine der Veränderlichen ${}^{\mathrm{m}}A$, ${}^{\mathrm{m}}C$, $A$, $B$, $C$, $A^{\mathrm{m}}$, $B^{\mathrm{m}}$ den Werth Null berührt oder einen negativen reellen Theil erhält. Die letzte jener Gleichungen, $n$ gleich Null gesetzt, gilt also sowohl für das System $(a,b,c)$ wie für $(\alpha,\beta,\gamma)$ und da
\[
\begin{gathered}\mathrm{M}(\alpha,\beta)=\tfrac12(1+\surd{-1})\mathrm{M}(a,b),\\[7pt]
\mathrm{M}(\alpha,\gamma)=\mathrm{M}(a,c)=\mathrm{M}(a,b)\end{gathered}
\]
ist, wird $\frac{\pi}{2}\surd{-1}=\log\surd{-1}$; es hat also $\pi$ die gebräuchliche Bedeutung als Verhältniss-Zahl der Länge des Umfangs eines Kreises zu dessen Durchmesser.

\par Mit Hülfe dieser Betrachtung lässt sich auch die Richtigkeit der folgenden von \textsc{Gauss} aufgezeichneten Bemerkung erweisen, der ich die hier benutzte Bezeichnungsweise zu Grunde lege und ein von \textsc{Gauss} berechnetes Beispiel folgen lasse:]
\par \textit{Die Arithmetisch-Geometrischen Mittel gestalten sich anders, wenn man für ein $b', b'', b'''\ .\ .$ den negativen Werth wählt: doch sind alle Resultate in folgender Form begriffen:}
\[
\frac{1}{\mathfrak{M}(a,b)}=\frac{1}{\mathrm{M}(a,b)}+\frac{4ik}{\mathrm{M}(a,c)}
\]
[wo $k$ eine ganze reelle Zahl bedeutet.]

\par \textit{Beispiel für einen imaginären Werth des A. G. Mittels}
\[
\begin{array}[t]{r@{\,=\,}r l r}
a&3.0000000&\log\ .\ .&0.4771213\\
b&1.0000000&&0.0000000\\[6pt]
a'&2.0000000&&0.3010300\\
b'&1.7320508&&0.2385606\\[6pt]
a''&1.8660254&&0.2709175\\
b''&1.8612098&&0.2697953\\[6pt]
a'''&1.8636176&&0.2703568\\
b'''&1.8636159&&0.2703564\\[6pt]
a''''&1.8636167&&0.2703566\\\hline
\end{array}
\]
\[
\begin{array}[t]{r@{\,=\,}l l r}
a&3.0000000&\log\ .\ .&0.4771213\\
c&2.8284270&&0.4515450\\[6pt]
\frac12\mathbin{.}{}'a&2.9142\mathrm{i}35&&0.4645214\\[6pt]
\frac12\mathbin{.}{}'c&2.9129510&&0.4643332\\[6pt]
\frac14\mathbin{.}{}''a&2.9135822&&0.4644273\\[4pt]\hline
\end{array}
\]
\SourceLeaf{388}
\begingroup\small\setlength{\arraycolsep}{2pt}
\[
\begin{array}{r@{\,=\,}r@{}l r r r r}
a&3.0000000&&\log\ .\ .\ 0.4771213&0&&\\
b&1.0000000&&0.0000000&360^{0}&&\\[6pt]
a'&2.0000000&&0.3010300&0&&\\
b'&-1.7320508&&0.2385606&180^{0}&&\\[6pt]
a''&0.1339746&&9.1270225&0&&\\
b''&&{}+1.8612098\,i&0.2697953&90^{0}&&\\[6pt]
a'''&0.0669873&{}+0.9306049\,i&9.9698876&85^{0}&52'&58''10\\
b'''&0.3530969&{}+0.3530969\,i&9.6984089&45&0&0\\[6pt]
a''''&0.2100421&{}+0.6418509\,i&9.8295254&71^{0}&52&46.58\\
b''''&0.2836930&{}+0.6208239\,i&9.8341482&65&26&29.05\\[6pt]
a^{\mathrm{V}}&0.2468676&{}+0.6313374\,i&9.8311572&68&38&36.05\\
b^{\mathrm{V}}&0.2470649&{}+0.6324002\,i&9.8318368&68&39&37.82\\[6pt]
a^{\mathrm{VI}}&0.24699625&{}+0.6318688\,i&9.8314971&68&39&6.95\\
b^{\mathrm{VI}}&0.24699625&{}+0.6318685\,i&9.8314970&68&39&6.93\\[6pt]
a^{\mathrm{VII}}&0.24699625&{}+0.63186865\,i&9.83149705&68&39&6.94\\\hline
\end{array}
\]
\endgroup
\[
\begin{aligned}&\frac{1}{\mathfrak{M}(a,b)}=+0.5365910-1.3728774\,i\\&=\frac{1}{\mathrm{M}(a,b)}+\frac{4i}{\mathrm{M}(a,c)}\end{aligned}
\]
\AGMArticle{13}
[Die Differentiale erster Ordnung der einzelnen Glieder des vollständigen Algorithmus eines arithmetisch geometrischen Mittels sind in Artikel 9 auf zweierlei Weise zu einander in Beziehung gesetzt. Die eine umfasst nur die Differentiale der Logarithmen der Quotienten jener Grössen, sie ergibt sich aus der wiederholten Anwendung der beiden Relationen $aa=bb+cc$, $\frac{a'}{c'}=\frac{a+b}{a-b}$, und kann durch Benutzung der im vorhergehenden Artikel gewonnenen Resultate zu einer Werthbestimmung des Differentials vom Quotienten zweier zusammengehöriger arithmetisch geometrischer Mittel erweitert und so dargestellt werden, dass $\frac{1}{cc}\,\mathrm{d}\log\frac{a}{b}$, welches mit $\Delta$ bezeichnet werden mag, für jedes positive und negative $n$, wenn nemlich ein negativer nachstehender Index $n$ als gleichbedeutend mit einem voranstehenden positiven Index von gleicher absoluter Grösse gedeutet wird,
\SourceLeaf{389}
\[
\begin{aligned}&\\&=\frac{1}{2^{n}}\mathbin{.}\frac{1}{a^{\mathrm{n}}a^{\mathrm{n}}}\,\mathrm{d}\log\frac{c^{\mathrm{n}}}{b^{\mathrm{n}}}\\&=\frac{1}{2^{n}}\mathbin{.}\frac{1}{b^{\mathrm{n}}b^{\mathrm{n}}}\,\mathrm{d}\log\frac{c^{\mathrm{n}}}{a^{\mathrm{n}}}\\&=\frac{1}{2^{n}}\mathbin{.}\frac{1}{c^{\mathrm{n}}c^{\mathrm{n}}}\,\mathrm{d}\log\frac{a^{\mathrm{n}}}{b^{\mathrm{n}}}\\&=\frac{\pi}{2}\mathbin{.}\frac{1}{\mathrm{M}(a,b)\mathbin{.}\mathrm{M}(a,c)}\,\mathrm{d}\log\frac{\mathrm{M}(a,c)}{\mathrm{M}(a,b)}\end{aligned}
\]
ist. Die andere Beziehung erstreckt sich auch mit auf die Differentiale der Logarithmen von zweigliedrigen Producten, sie folgt aus $b'b'=ab$ und ${}'c{}'c=4ac$ in der Form
\[
\begin{aligned}
\mathrm{d}\log(a^{\mathrm{m}}\mathbin{.} b^{\mathrm{m}})&=\mathrm{d}\log(a^{\mathrm{m}+1}\mathbin{.} b^{\mathrm{m}+1})\\&\quad-\Delta\mathbin{.}2^{m+1}\mathbin{.} c^{\mathrm{m}+1}\mathbin{.} c^{\mathrm{m}+1}\\[3pt]
\mathrm{d}\log({}^{\mathrm{m}}a\mathbin{.}{}^{\mathrm{m}}c)&=\mathrm{d}\log({}^{\mathrm{m}+1}a\mathbin{.}{}^{\mathrm{m}+1}c)\\&\quad+\Delta\mathbin{.}\frac{1}{2^{m+1}}\mathbin{.}{}^{\mathrm{m}+1}b\mathbin{.}{}^{\mathrm{m}+1}b
\end{aligned}
\]
und ergibt, wenn man $m$ bis zur unendlichen Grenze wachsen lässt, die in Artikel 9 aufgestellte Gleichung und die dieser entsprechende nemlich:
\[
\begin{aligned}
2\,\mathrm{d}\log\mathrm{M}(a,b)&=\mathrm{d}\log(a^{\mathrm{n}}\mathbin{.} b^{\mathrm{n}})\\&\quad+\Delta\{2^{n+1}\mathbin{.} c^{\mathrm{n}+1}\mathbin{.} c^{\mathrm{n}+1}+\ .\ +2^{n+m}\mathbin{.} c^{\mathrm{n}+\mathrm{m}}\mathbin{.} c^{\mathrm{n}+\mathrm{m}}+\ .\}\\[3pt]
2\,\mathrm{d}\log\mathrm{M}(a,c)&=\mathrm{d}\log({}^{\mathrm{n}}a\mathbin{.}{}^{\mathrm{n}}c)\\&\quad-\Delta\{2^{-n-1}\mathbin{.}{}^{\mathrm{n}+1}b\mathbin{.}{}^{\mathrm{n}+1}b+\ .\ +2^{-n-m}\mathbin{.}{}^{\mathrm{n}+\mathrm{m}}b\mathbin{.}{}^{\mathrm{n}+\mathrm{m}}b+\ .\}
\end{aligned}
\]
\par Hieraus kann durch Elimination der Differentiale mit Hülfe des oben gefundenen Ausdrucks für $\Delta$ die Gleichung
\[
\begin{aligned}
\frac{4}{\pi}\mathrm{M}(a,b)\,\mathrm{M}(a,c)&=.\ .-2^{-\mu+n}\mathbin{.}{}^{\mu-\mathrm{n}}b\mathbin{.}{}^{\mu-\mathrm{n}}b\\&\quad-.\ .-2^{n-1}\mathbin{.} b^{\mathrm{n}-1}\mathbin{.} b^{\mathrm{n}-1}\\
&{}+2^{n}\mathbin{.} a^{\mathrm{n}}\mathbin{.} a^{\mathrm{n}}-2^{n+1}\mathbin{.} c^{\mathrm{n}+1}\mathbin{.} c^{\mathrm{n}+1}\\&\quad-.\ .-2^{n+m}\mathbin{.} c^{\mathrm{n}+\mathrm{m}}\mathbin{.} c^{\mathrm{n}+\mathrm{m}}-.\ .
\end{aligned}
\]
abgeleitet werden, welche \textsc{Gauss} neben der im vorigen Artikel wiedergegebenen ersten Gleichung für $\frac{\pi}{2}\frac{\mathrm{M}(a,b)}{\mathrm{M}(a,c)}$ aufgezeichnet hat, ohne den Weg, auf welchem sie gefunden waren, anzudeuten.]
\AGMArticle{14}
\par [Für die vollständigen Differentiale zweiter Ordnung findet man unmittelbar aus den Ausdrücken für $\Delta$, dass jedes mit gemeinsamem Index behaftete System von Gliedern $a,b,c$ die drei Gleichungen
\[
\begin{aligned}
\frac{1}{\Delta}\,\mathrm{d}\log\frac{c}{b}\mathbin{.}\mathrm{d}\log a&=\tfrac12\,\mathrm{d}\left(\frac{1}{\Delta}\,\mathrm{d}\log\frac{c}{b}\right)\\[4pt]
\frac{1}{\Delta}\,\mathrm{d}\log\frac{c}{a}\mathbin{.}\mathrm{d}\log b&=\tfrac12\,\mathrm{d}\left(\frac{1}{\Delta}\,\mathrm{d}\log\frac{c}{a}\right)\\[4pt]
\frac{1}{\Delta}\,\mathrm{d}\log\frac{a}{b}\mathbin{.}\mathrm{d}\log c&=\tfrac12\,\mathrm{d}\left(\frac{1}{\Delta}\,\mathrm{d}\log\frac{a}{b}\right)
\end{aligned}
\]
 erfüllt. Multiplicirt man darin die Zähler und Nenner unter den zweimal zu differentiirenden Logarithmen der Reihe nach mit $a,b,c$ und löst alle Logarith\SourceLeaf{390}men
von Quotienten in Differenzen von Logarithmen auf, so ersieht man leicht,
dass der Ausdruck
\[
\frac{1}{\Delta}\,\mathrm d\log a^{\mathrm n}\mathbin{.}\mathrm d\log b^{\mathrm n}
-\frac12\mathrm d\left(\frac{1}{\Delta}\,\mathrm d\log a^{\mathrm n}b^{\mathrm n}\right)
\]
der mit $D$ bezeichnet werden soll, seinen Werth nicht ändert, wenn man statt
$a^{\mathrm n}$ und $b^{\mathrm n}$ setzt: $a^{\mathrm n}$ und $c^{\mathrm n}$ oder:
$b^{\mathrm n}$ und $c^{\mathrm n}$. Nimmt man das Mittel der beiden letzten so
entstandenen Ausdrücke und berücksichtigt, dass
\[
\mathrm d\log a^{\mathrm n}+\mathrm d\log b^{\mathrm n}
=2\mathrm d\log b^{\mathrm n+1}
\]
ist, so folgt, dass man in jenem Ausdruck ohne dessen Werth zu ändern statt
der beiden genannten Grössen auch $c^{\mathrm n}$ und $b^{\mathrm n+1}$ setzen kann
und ferner, wenn man in diesem wieder $\mathrm d\log c^{\mathrm n}$ durch
$\frac12\mathrm d\log a^{\mathrm n+1}+\frac12\mathrm d\log c^{\mathrm n+1}$ ersetzt,
dass man mit demselben Erfolge $a^{\mathrm n+1}$ und $b^{\mathrm n+1}$ statt
$a^{\mathrm n}$ und $b^{\mathrm n}$ setzen kann. Es ist also der Werth jenes
Differential-Ausdrucks unabhängig von $n$ und demnach gleich
\[
\begin{aligned}
D&=\displaystyle\frac{1}{\Delta}\,\mathrm d\log b^{\mathrm n}\mathbin{.}\mathrm d\log c^{\mathrm n}
-\frac12\mathrm d\left(\frac{1}{\Delta}\,\mathrm d\log b^{\mathrm n}c^{\mathrm n}\right)\\[6pt]
&=\displaystyle\frac{1}{\Delta}\,\mathrm d\log c^{\mathrm n}\mathbin{.}\mathrm d\log a^{\mathrm n}
-\frac12\mathrm d\left(\frac{1}{\Delta}\,\mathrm d\log c^{\mathrm n}a^{\mathrm n}\right)\\[6pt]
&=\displaystyle\frac{1}{\Delta}\,\mathrm d\log a^{\mathrm n}\mathbin{.}\mathrm d\log b^{\mathrm n}
-\frac12\mathrm d\left(\frac{1}{\Delta}\,\mathrm d\log a^{\mathrm n}b^{\mathrm n}\right)\\[6pt]
&=\displaystyle\mathrm M(a,b)\,\mathrm d\left\{\frac{1}{\Delta}\,\mathrm d\frac{1}{\mathrm M(a,b)}\right\}\\&=\mathrm M(a,c)\,\mathrm d\left\{\frac{1}{\Delta}\,\mathrm d\frac{1}{\mathrm M(a,c)}\right\}
\end{aligned}
\]
\par Setzt man $a$ als unveränderlich voraus, so wird
\[
D=\Delta bbcc=-bb\,\mathrm d\log b=cc\,\mathrm d\log c
\]
und es entsteht die in Artikel 8 aus der Reihenentwickelung abgeleitete für
$\frac{a}{\mathrm M(a,c)}$ und $\frac{a}{\mathrm M(a,b)}$ als Werthe von $\mu$
geltende Differentialgleichung
\[
\mathrm{dd}\mu-\frac{\mathrm d\Delta}{\Delta}\mathbin{.}\mathrm d\mu
-\Delta\Delta bbcc\mathbin{.}\mu=0
\]
aus welcher sich nach den Untersuchungen in der Abhandlung
`\textit{Determinatio seriei nostrae per aequationem differentialem secundi ordinis}'
auch wieder die Darstellung durch die \textsc{Gauss}ischen Reihen:
\[
\begin{gathered}\frac{a}{\mathrm M(a,b)}=F\left(\frac12,\frac12,1,\frac{cc}{aa}\right),\\[7pt]
\frac{a}{\mathrm M(a,c)}=F\left(\frac12,\frac12,1,\frac{bb}{aa}\right)\end{gathered}
\]
\SourceLeaf{391}
und als specielle Fälle der dortigen Gleichungen [90] und [96] der oben Art. 12
gefundene Grenzwerth von $\mathrm M(1,\varepsilon)\log\frac{4}{\varepsilon}$ und die
im vorigen Artikel aufgestellte Differentialgleichung erster Ordnung zwischen
$\mathrm M(a,b)$ und $\mathrm M(a,c)$ ergeben.]

\AGMArticle{15}
[Die Differentialgleichung zweiter Ordnung für die Glieder der Reihe des
arithmetisch geometrischen Mittels nimmt eine Form an, in welcher sie
Differentiale nur von Quotienten der Veränderlichen enthält, wenn man mit einer
beliebigen Grösse $e$ den Ausdruck
$D+\mathrm d(\frac{1}{\Delta}\,\allowbreak \mathrm d\log e)+\frac{1}{\Delta}(\mathrm d\log e)^2$
bildet, dieser wird nemlich
\[
\begin{aligned}&-\surd(eea^{\mathrm n}b^{\mathrm n})\mathbin{.}\mathrm d
\left(\frac{1}{\Delta a^{\mathrm n}b^{\mathrm n}}\mathbin{.}\mathrm d
\surd\frac{a^{\mathrm n}b^{\mathrm n}}{ee}\right)\\&\quad-\frac{1}{4\Delta}\left(\mathrm d\log\frac{a^{\mathrm n}}{b^{\mathrm n}}\right)^2\end{aligned}
\]
ein Ausdruck, dessen Werth also unabhängig von $n$ ist und sich auch nicht
ändert, wenn man $b^{\mathrm n}$ und $c^{\mathrm n}$ oder $c^{\mathrm n}$ und
$a^{\mathrm n}$ statt $a^{\mathrm n}$ und $b^{\mathrm n}$ setzt. Derselbe verwandelt
sich für beständig wachsende positive und für negative $n$ in
\[
\begin{gathered}-e\mathrm M(a,b)\mathbin{.}\mathrm d
\left\{\frac{1}{\Delta\mathrm M(a,b)^2}\,\mathrm d\frac{\mathrm M(a,b)}{e}\right\}
\\[7pt]\mbox{und in}\\[7pt]
-e\mathrm M(a,c)\mathbin{.}\mathrm d
\left\{\frac{1}{\Delta\mathrm M(a,c)^2}\,\mathrm d\frac{\mathrm M(a,c)}{e}\right\}\end{gathered}
\]
dagegen für $n$ gleich Null und für $a$, $b$, $c$ als besondere Werthe von $e$
beziehungsweise in $+\Delta bbcc,\quad-\Delta ccaa,\quad-\Delta aabb.$

\par Setzt man also zur Abkürzung
\[
\begin{gathered}\surd\frac{a}{\mathrm M(a,b)}=p,\\[7pt]
\surd\frac{b}{\mathrm M(a,b)}=q,\\[7pt]
\surd\frac{c}{\mathrm M(a,b)}=r,\\[7pt]
-\pi\frac{\mathrm M(a,b)}{\mathrm M(a,c)}=\log y\end{gathered}
\]
so wird mit Rücksicht auf die Gleichung für $\Delta$ in Art. 13:
\[
\begin{aligned}
\displaystyle\frac{\Delta}{2}\mathrm M(a,b)^2
&=\displaystyle\frac14\mathrm d\log y\\&=\frac{1}{p^4}\,\mathrm d\log\frac{r}{q}\\&=\frac{1}{q^4}\,\mathrm d\log\frac{r}{p}\\&=\frac{1}{r^4}\,\mathrm d\log\frac{p}{q}\\[8pt]
&=\displaystyle-\frac{pp}{q^4r^4}\,\mathrm d
\left(\frac{1}{\mathrm d\log y}\mathbin{.}\mathrm d\frac{1}{pp}\right)\\&=\frac{qq}{r^4p^4}\,\mathrm d
\left(\frac{1}{\mathrm d\log y}\mathbin{.}\mathrm d\frac{1}{qq}\right)\\&=\frac{rr}{p^4q^4}\,\mathrm d
\left(\frac{1}{\mathrm d\log y}\mathbin{.}\mathrm d\frac{1}{rr}\right)
\end{aligned}
\]
und von gleicher Form werden die Ausdrücke des $-\frac{\Delta}{2}\mathrm M(a,c)^2$ in
\[
\begin{gathered}\surd\frac{a}{\mathrm M(a,c)},\\[7pt]
\surd\frac{c}{\mathrm M(a,c)},\\[7pt]
\surd\frac{b}{\mathrm M(a,c)},\\[7pt]
-\pi\frac{\mathrm M(a,c)}{\mathrm M(a,b)}\end{gathered}
\]
\par Die Elimination von je zwei der drei Grössen $p,q,r$ ergibt
\[
\begin{aligned}&\left\{\frac{1}{p^4}\frac{1}{\mathrm d\log y}\,\mathrm d\log
\left[\frac{16}{p^6}\frac{1}{\mathrm d\log y}\,\mathrm d
\left(\frac{1}{\mathrm d\log y}\mathbin{.}\mathrm d\frac{1}{pp}\right)\right]\right\}^{2}\\&\quad-\frac{16}{p^6}\frac{1}{\mathrm d\log y}\,\mathrm d
\left(\frac{1}{\mathrm d\log y}\mathbin{.}\mathrm d\frac{1}{pp}\right)-1=0\end{aligned}
\]
